\subsection{The alternate source datum over \(\mathbf Q(\sqrt{-7})\)}
\label{sec:minus7-source-coverage}

\begin{lemma}[K17-MINUS7-SOURCE]\label{lem:minus7-source}
Put
\[
 \ell_0=25078104767246023,\qquad E=\mathbf Q(\sqrt{-7}),
 \qquad e=16,\qquad N=2,\qquad\psi=1.
\]
The actual source datum \(\mathrm{GD}(\Omega,E,\ell_0,16,1)\) holds.
For this datum, the direct criterion of
Lemma~\ref{lem:step17-unit} may be applied at a place of \(E\) above
the split prime \(2\).
\end{lemma}

\begin{proof}
The definition \texttt{HeightAndVanishing/Counts.lean:GD} in
\cite{K17PinnedLeanMirror} requires an odd coefficient prime, positive
\(e\), even \(N\), the prime-to-\(\ell_0\) and quadratic conditions on
\(\psi\), the parity condition \(\psi(-1)=(-1)^e\), an imaginary
quadratic field, coprimality of its discriminant with \(N\), and
\(\ell_0\nmid |d_E|h_E\). Here \(d_E=-7\) and \(h_E=1\): Minkowski's
ideal-class bound is \((2/\pi)\sqrt7<2\), so every class has an integral
ideal of norm one. All these conditions follow at once. Also
\(\ell_0-1>32=2e\), which gives the source's condition
\(\ell_0-1\nmid2e\).

Its remaining condition \texttt{eta}, after \(\psi=1\), says that the
kernel of the mod-\(\ell_0\) cyclotomic character is not contained in
\(G_E\). Equivalently, \(E\not\subset\mathbf Q(\mu_{\ell_0})\).
The latter extension is unramified at \(7\), whereas \(E\) is ramified
there. This proves the condition. It is the same ramification argument
as the source's \texttt{SmallClassNumbers/ClassNumberThree.lean:gd},
with the ramified prime \(7\) in place of \(3\).

The polynomial \(X^2-X+2\) for
\((1+\sqrt{-7})/2\) has two distinct roots modulo \(2\), so \(2\)
splits in \(E\). The exact statement \texttt{OneCharU\_s} in
\texttt{DenseThresholdCases/NumberU3US.lean} permits every
\(s\in G_E\) and every specified comparison trace. It places no lower
bound on an auxiliary prime and no condition that it be prime to \(N\).
The restrictions \(p\ge7\) and \(p\nmid N\) occur in the separate
counting statement \texttt{NumberU3U\_s}, which is not used here.
Taking \(s\) to be geometric Frobenius above \(2\) is therefore allowed.
\end{proof}

\begin{lemma}[K17-MINUS7-U3-COVERAGE]\label{lem:minus7-u3-coverage}
Let \(G=U(I_3/E)\), and let
\[
 K=\prod_v\operatorname{Stab}_{G(\mathbf Q_v)}
       ((O_E\otimes\mathbf Z_v)^3).
\]
Every actual source representation
\(\rho\in\mathrm{Adm3U\_s}(\iota,16,2)\) has its split-prime Hecke
system in the full algebraic modular-form space for \(G,K\) and
compact highest weight \((15,0,-15)\).
If \(\alpha\) is its eigenvalue for \(2^{15}T_2\), with \(T_2\) the
positive minuscule operator, then
\[
 \iota\operatorname{tr}\rho(s)=\frac{\alpha}{2^{16}},
 \qquad s=\operatorname{Frob}_2^{-1}.
 \tag{S7.1}
\]
\end{lemma}

\begin{proof}
The source-to-automorphic part of
Corollary~\ref{cor:u3-general-coverage} does not use the particular
imaginary quadratic field. We give the actual substitutions.
\texttt{Adm3U\_s} retains \texttt{Adm3Alt} and \texttt{Th.Unl};
it supplies the literal condition
\(\mathrm{Level}(\ell_0,2,(16,0,-16),\rho)\) and
\(\mathrm{IsAutCsdA}(\iota,r)\) for \(r=\rho\chi_{\ell_0}^{-1}\).
Every \texttt{KillsMu 2} character is trivial. Reusing the same level
witnesses and applying \texttt{level\_twist\_cyclotomic\_pow} gives
\[
 \mathrm{Level}(\ell_0,1,(15,-1,-17),r).
 \tag{S7.2}
\]
This gives continuity and hence the actual \texttt{WitA} witness.
The proved \texttt{unlAll\_of} and the common cyclotomic twist give
the above-\(\ell_0\) \texttt{Unl1} hypotheses for \(r\).
Apply \texttt{Keep.levelKA\_alt} at modulus one, with all these
hypotheses retained. Its full conclusion is a regular algebraic,
conjugate-self-dual cuspidal representation \(\Pi\) of
\(\mathrm{GL}_3(\mathbf A_E)\), unramified at every finite place,
including above \(\ell_0\), with
\[
 \operatorname{InfChar}(\Pi)=\operatorname{infOf}(15,-1,-17)
 =(16,0,-16),
\]
and exact Frobenius--Satake compatibility away from \(\ell_0\).
The unitary infinity parameter is
\(\operatorname{diag}((z/\bar z)^{16},1,(z/\bar z)^{-16})\).
Thus its compact highest weight is
\((16,0,-16)-(1,0,-1)=(15,0,-15)\).
The source declaration and normalization locators are those in the
complete proof of Lemma~\ref{lem:step17-u3-level}.

We check every change in the compact transfer argument of
Theorem~\ref{thm:u3-coverage}. The group \(U(I_3/E)\) is a pure inner
form of the quasi-split unitary group and is quasi-split at all finite
places. The three distinct infinity characters are
conjugate-orthogonal. Thus \(\Pi\) gives a simple generic parameter
for the untwisted base-change embedding. Its reduced component group
is trivial, so the proved generic pure-inner multiplicity formula
\cite[Theorem~5.0.5, p.~206; \S1.3.4--1.3.5, pp.~69--70]{K17KMSW2014}
permits any relevant local choices. Shin's odd-rank theorem gives
temperedness at every finite place
\cite[Theorem~1.2 and Corollary~1.3, pp.~1647--1648;
Corollary~7.9, p.~1728]{K17Shin2011}.
The central character of \(\Pi\) is unramified everywhere finite and
trivial at infinity; it factors through the trivial ideal class group
of \(E\), and hence is trivial.

At split places the usual spherical member supplies the required
fixed vector. At inert unramified places the standard lattice
stabilizer is hyperspecial, and
\cite[Proposition~1.5.2(5), p.~79]{K17KMSW2014} supplies the
unramified packet member. This covers the dyadic place as well,
since \(2\) is split in the present field.

It remains to check the ramified prime \(7\) and its full lattice
stabilizer. Over \(\mathbf Q_7\), Hensel's lemma gives a unit
\(a\in\mathbf Z_7\) with \(a^2=-5\), starting from \(a=3\bmod7\).
Then \(v=(1,2,a)\) is a primitive isotropic vector for \(I_3\).
Choose an integral \(w\) with \(h(v,w)=1\); replacing \(w\) by
\(w-h(w,w)v/2\) gives an integral hyperbolic plane. Its orthogonal
complement has unit norm \(d\). For the resulting integral basis
matrix \(B\), determinants give
\(-d=N_{E_7/\mathbf Q_7}(\det B)\); scaling the remaining basis
vector by \((\det B)^{-1}\) makes its norm \(-1\). This identifies
the lattice integrally with the Adler--Lansky model
\[
 \Phi=\begin{pmatrix}0&0&1\\0&-1&0\\1&0&0\end{pmatrix}.
\]
The full stabilizer has reductive quotient
\(O_3(\mathbf F_7)=SO_3(\mathbf F_7)\times\{\pm1\}\);
it is the special maximal subgroup at their vertex \(y\)
\cite[\S4.1--4.2.1, pp.~10--12]{K17AdlerLansky2010}.

Conjugate self-duality and trivial determinant put the unramified
base-change roots at \(7\) in the form \(\{u,1,u^{-1}\}\), with \(u\)
unitary. Induce from the unitary Borel character
\(\lambda_1=\chi,\lambda_2=1\), where \(\chi\) is unramified with
Frobenius value \(u\). The lifted \(\mathrm{GL}_3(E_7)\) principal
series is irreducible, because its unitary inducing characters have
no ratio \(|\cdot|^{\pm1}\). The unitary principal series is either
irreducible or is the two-constituent packet for nontrivial
unramified quadratic \(\chi\).
The exact reducibility and base-change assertions are
\cite[\S4.4, pp.~14--15; Proposition~6.1(1),(4), p.~21;
Proposition~7.1(i),(iv), p.~28]{K17AdlerLansky2010}.
Their needed principal-series statements follow directly from
\cite[Proposition~5.1 and its complete proof, pp.~19--20,
published numbering Proposition~4.1]{K17AdlerLanskyUnram},
using Rogawski's Proposition~4.10.2 and the paragraph before
Theorem~13.2.1. Thus neither the separate depth-preservation
Hypothesis~1.5 nor the \(q\ne7\) restriction in the unused
supercuspidal theorem enters this argument.
Iwasawa decomposition gives a one-dimensional fixed space in the
induction for the full stabilizer. Exactness of compact-open
invariants in characteristic zero gives a fixed vector in one of
the constituents, all of which lie in the required packet.

Choose these finite members and the compact infinity member.
The cited global multiplicity formula gives multiplicity one, and
the algebraic modular-form description places the system in the
full space asserted. Finally, the source's cyclotomic twist and
\(\operatorname{algSat}(3,2)\) both contribute a factor \(2\), which
cancel. The positive minuscule Satake formula is
\(T_2=2\sum\beta\). Hence \(2^{15}T_2=2^{16}\sum\beta\), proving
(S7.1). An inverse-minuscule finite model must use its separately
verified conjugation identification with this positive operator.
\end{proof}

\begin{lemma}[K17-MINUS7-GL2-CONDUCTOR]\label{lem:minus7-gl2-conductor}
Every \(\rho\in\mathrm{Adm2A}(\iota,16,14)\) has, at the same
geometric Frobenius \(s\), the trace
\[
 \iota\operatorname{tr}\rho(s)
   =1+\frac{\zeta\,a_2(f)}{2^{16}},
 \qquad f\in S_{33}(\Gamma_1(7)),\qquad \zeta^3=1,
 \tag{S7.3}
\]
where \(f\) is a normalized primitive cuspidal eigenform of
conductor exactly \(7\). The original untwisted eigenform need
not have conductor \(7\).
\end{lemma}

\begin{proof}
The exact definition in \texttt{AutNewformW.lean} supplies a
continuous irreducible \(W\), a weight-\(33\) newform attached to
\(R=W\chi_{\ell_0}^{-16}\), the literal
\(\mathrm{Level}(\ell_0,14,(16,-16),W)\), and
\(\rho\simeq1\oplus W|_{G_E}\). We retain its other
\(\mathrm{WitAQ}\) and determinant conditions.
Attachment at almost all primes identifies \(R\) with the
Deligne representation of the primitive eigenform, by
Chebotarev and Brauer--Nesbitt, as in
Lemma~\ref{lem:gl2-conductor}.

Since \(\mathbf Q(\mu_{14})=\mathbf Q(\mu_7)\), every source
\texttt{KillsMu 14} character is a global Dirichlet character
modulo \(7\). Away from \(\ell_0\), the literal inertia
diagonalization therefore gives finite inertia, zero monodromy,
and no ramification outside \(7\). At \(7\) its two diagonal
characters are the restrictions of some global mod-\(7\)
characters \(\kappa_1,\kappa_2\). Both may be nontrivial.
We do not discard that possibility.

At the coefficient prime, the ordinary flag for \(R\) has
unramified characters times the cyclotomic powers \(1\) and
\(\chi_{\ell_0}^{-32}\). These unramified characters have unit
Frobenius values. Saito's classical full Weil--Deligne
compatibility theorem
\cite[\S2c, Theorem~3, p.~9]{K17SaitoSurvey}, and the exactness
of \(D_{\mathrm{st}}\)
\cite[\S4.1, Proposition~4.2, p.~23]{K17ColmezFontaine2000},
give two Frobenius slopes differing by \(32\).
The relation \(N\varphi=\ell_0\varphi N\) would force difference
one if \(N\ne0\). Thus \(N=0\). The finite Weil--Deligne
inertia action has trivial graded characters and hence is
trivial by characteristic-zero semisimplicity. This proves
unramifiedness at \(\ell_0\), with exactly the coefficient-prime
argument already detailed in Lemma~\ref{lem:gl2-conductor}.
At the other primes use
\cite[\S0.7, Theorem~A, pp.~410--411]{K17Carayol1986}.

Twist the automorphic representation and \(R\) by the global
finite character \(\kappa_1^{-1}\). Twisting is multiplication
of an adelic cusp form by \(\kappa_1^{-1}(\det g)\); it preserves
cuspidality and the weight-\(33\) holomorphic discrete series.
It multiplies every good-prime Hecke eigenvalue by the same
Dirichlet value, and its Galois realization is the indicated
tensor twist. The new inertia at \(7\) is
\(1\oplus\kappa_2\kappa_1^{-1}\).
The original nebentypus is odd, because the weight is odd.
It is unramified outside \(7\), and twisting its determinant
by a square preserves oddness. Consequently
\(\kappa_2\kappa_1^{-1}\) is a nontrivial odd mod-\(7\)
character. Zero monodromy and this tame inertia give conductor
exponent exactly one at \(7\). Everywhere else the twisted
representation is unramified. Conductor-preserving local
compatibility therefore puts its primitive form \(f\) in
\(S_{33}(\Gamma_1(7))\), with conductor exactly \(7\).

Untwisting multiplies \(a_2(f)\) by a value of a mod-\(7\)
Dirichlet character at \(2\). Since \(2\) has order three in
\((\mathbf Z/7)^\times\), that value belongs to \(\mu_3\).
The source twist gives
\(R(s)=2^{16}W(s)\), and adjoining the trivial summand yields
(S7.3). Inverting a reciprocity convention only permutes
the same three allowed roots of unity.
\end{proof}

\begin{corollary}[K17-MINUS7-UNIT]\label{cor:minus7-unit}
Put \(E_2=1+2^{16}+2^{32}\) and \(t_2=1+2^{32}\).
Work over the valuation ring \(\mathcal O\) of a finite
coefficient extension of \(\mathbf Q_{\ell_0}\), large enough
for both coefficient representations and for \(\mu_3\).
Suppose the entire compact space of
Lemma~\ref{lem:minus7-u3-coverage} has an
\(\mathcal O\)-integral Hecke lattice and
\[
 \det(E_2I-2^{15}T_2)\in\mathcal O^{\times}.
 \tag{S7.4}
\]
Suppose also that the full space \(S_{33}(\Gamma_1(7))\)
has an \(\mathcal O\)-integral Hecke lattice with
\[
 \det(t_2^3 I-T_2^3)\in\mathcal O^{\times}.
 \tag{S7.5}
\]
Then the source Selmer length for \(e=16,\psi=1\) is zero.
\end{corollary}

\begin{proof}
Every admitted irreducible trace difference is
\((\alpha-E_2)/2^{16}\) by (S7.1).
The roots of the integral compact characteristic polynomial
are integral; (S7.4) therefore makes every such factor a unit.
Every admitted rank-two difference is
\((\zeta a_2(f)-t_2)/2^{16}\) by (S7.3).
Over a coefficient extension containing \(\mu_3\),
\[
 \det(t_2^3I-T_2^3)
   =\prod_{\zeta^3=1}\det(t_2I-\zeta T_2).
\]
All eigenvalues are integral, so (S7.5) makes each of their
individual differences a unit at every prime above \(\ell_0\).
The denominator is a unit. Lemma~\ref{lem:minus7-source}
supplies the exact datum and permits this Frobenius.
Lemma~\ref{lem:step17-unit} now forces the length to vanish.
This uses pointwise coverage of every actual source predicate;
no replacement of the source's \texttt{TrData} or comparison
of its multiplicities is asserted.
\end{proof}
